\documentclass[a4paper,10pt]{article}
\usepackage{lamia-pedagogy}

\begin{document}
% PAGE 3 : Compétence 02 - Lecture
% =========================================================================
\ficheHeader{1 Heure}{\makebox[2.5cm]{\dotfill}}{1\textsuperscript{ère} AM}{2026 / 2027}

\vspace{-2pt}
\competenceBanner{\faBookOpen\ Compétence 02 : Lecture et phonétique}

\begin{minipage}[t]{0.48\textwidth}
\begin{tcolorbox}[
  enhanced,
  colback=secondarylight,
  colframe=secondary,
  arc=1.5mm,leftrule=3pt,rightrule=0.5pt,toprule=0.5pt,bottomrule=0.5pt,
  top=3pt,bottom=3pt,left=6pt,right=6pt,
  title={\small\textbf{\faBullseye\ Critères \& Objectifs}},
  coltitle=secondary!90!black,attach boxed title to top left={yshift=-2mm,xshift=4mm},
  boxed title style={colback=white,boxrule=0.5pt,colframe=secondary,arc=1mm,top=1pt,bottom=1pt}
]
\vspace{1pt}
{\scriptsize
\begin{itemize}[leftmargin=*,noitemsep,topsep=1pt]
  \item[\faCheck] Activer la correspondance graphie/phonie.
  \item[\faCheck] Réaliser une lecture fluide et expressive.
  \item[\faCheck] Respecter l'intonation des différents types de phrases.
  \item[\faCheck] Distinguer les sons : \textbf{/g/ -- /3/}, \textbf{/s/ -- /z/}, \textbf{/k/ -- /s/}.
\end{itemize}
}
\end{tcolorbox}
\end{minipage}
\hfill
\begin{minipage}[t]{0.48\textwidth}
\begin{tcolorbox}[
  enhanced,
  colback=accentlight,
  colframe=accent,
  arc=1.5mm,leftrule=3pt,rightrule=0.5pt,toprule=0.5pt,bottomrule=0.5pt,
  top=3pt,bottom=3pt,left=6pt,right=6pt,
  title={\small\textbf{\faExclamationTriangle\ Obstacles \& Difficultés}},
  coltitle=accent!90!black,attach boxed title to top left={yshift=-2mm,xshift=4mm},
  boxed title style={colback=white,boxrule=0.5pt,colframe=accent,arc=1mm,top=1pt,bottom=1pt}
]
\vspace{1pt}
{\scriptsize
\begin{itemize}[leftmargin=*,noitemsep,topsep=1pt]
  \item[\faTimesCircle] Difficulté à associer la lettre au phonème correspondant.
  \item[\faTimesCircle] Décodage laborieux gênant la fluence.
  \item[\faTimesCircle] Difficulté à moduler l'intonation (expressivité).
\end{itemize}
}
\end{tcolorbox}
\end{minipage}

\vspace{1pt}
\begin{tcolorbox}[
  enhanced,
  colback=secondarylight,
  colframe=secondary,
  arc=1.5mm,leftrule=3.5pt,rightrule=0.6pt,toprule=0.6pt,bottomrule=0.6pt,
  top=2.5pt,bottom=2.5pt,left=8pt,right=8pt
]
{\small
\textbf{\faUsers\ Modalité :} Travail en groupes \hfill \textbf{\faChalkboard\ Support :} Tableau, cahier
}
\end{tcolorbox}

\vspace{1pt}
\begin{tcolorbox}[
  enhanced,
  colback=white,
  colframe=cardborder,
  arc=1.5mm,
  top=2pt,bottom=2pt,left=6pt,right=6pt,
  title={\scriptsize\textbf{\faUsers\ Élèves ayant des difficultés de lecture}},
  coltitle=primary,attach boxed title to top left={yshift=-1.8mm,xshift=3mm},
  boxed title style={colback=primarylight,boxrule=0.4pt,colframe=cardborder,arc=1mm,top=0.5pt,bottom=0.5pt}
]
\vspace{1pt}
{\scriptsize
\begin{tabularx}{\linewidth}{@{}l X l X@{}}
1. & \makebox[\linewidth]{\dotfill} & 6. & \makebox[\linewidth]{\dotfill} \\
2. & \makebox[\linewidth]{\dotfill} & 7. & \makebox[\linewidth]{\dotfill} \\
3. & \makebox[\linewidth]{\dotfill} & 8. & \makebox[\linewidth]{\dotfill} \\
4. & \makebox[\linewidth]{\dotfill} & 9. & \makebox[\linewidth]{\dotfill} \\
5. & \makebox[\linewidth]{\dotfill} & 10. & \makebox[\linewidth]{\dotfill} \\
\end{tabularx}
}
\end{tcolorbox}

\vspace{1pt}
\begin{tcolorbox}[
  enhanced,
  colback=primarylight,
  colframe=primary,
  arc=1.5mm,
  top=2.5pt,bottom=2.5pt,left=6pt,right=6pt,
  title={\small\textbf{\faChalkboardTeacher\ Mise en {\oe}uvre de la remédiation}},
  coltitle=white,colbacktitle=primary
]
\vspace{1pt}
{\small
\begin{tabularx}{\linewidth}{|X|X|}
\hline
\rowcolor{secondarylight}
\textbf{\faUserTie\ Rôle de l'enseignante} & \textbf{\faUserGraduate\ Rôle des apprenants} \\
\hline
$\bullet$ Présente les activités et le matériel. & $\bullet$ Lisent silencieusement puis à haute voix. \\
$\bullet$ Explicite clairement les consignes. & $\bullet$ Écoutent et posent des questions. \\
$\bullet$ Oriente, régule et encourage. & $\bullet$ S'entraident et échangent en binômes. \\
\hline
\end{tabularx}
}
\end{tcolorbox}

\vspace{1pt}
\begin{tcolorbox}[
  enhanced,
  colback=accentlight,
  colframe=accent,
  arc=1.5mm,
  top=2.5pt,bottom=2.5pt,left=7pt,right=7pt
]
\textbf{\faLightbulb\ Moment de mise en confiance -- Devinette : Qui suis-je ?}\\[1pt]
{\small\itshape
\guillemotleft\ Je travaille à l'hôpital, je soigne les malades, j'utilise un stéthoscope. Qui suis-je ?\ \guillemotright\ \hfill \textbf{Réponse :} \makebox[3cm]{\dotfill}
}
\end{tcolorbox}

\vspace{1pt}
{\small
\badge[primary]{\faTasks\ Activités d'entraînement phonétique et de lecture}

\vspace{2pt}
\textbf{Activité 01 (Fluence de lecture) :} \textit{Lis attentivement le passage suivant :}\\
\begin{tcolorbox}[enhanced,colback=white,colframe=cardborder,arc=1mm,top=3pt,bottom=3pt,left=6pt,right=6pt]
\textit{\guillemotleft\ Chaque jour, Amira va à l'école à 8 heures du matin. Dans la classe, elle retrouve ses camarades et écoute sa maîtresse. Elle apprend à lire et à écrire, elle aime dessiner les fleurs et les animaux.\ \guillemotright}
\end{tcolorbox}

\vspace{2pt}
\textbf{Activité 02 (Discrimination de la lettre \guillemotleft\ C \guillemotright) :} \textit{Classe ces mots dans le tableau :}\\
\centerline{\textbf{Carotte -- cerise -- coq -- bicyclette -- cube -- citron}}
\vspace{2pt}
\begin{tabularx}{\linewidth}{|X|X|}
\hline
\rowcolor{secondarylight}
\centering\textbf{La lettre \guillemotleft\ C \guillemotright\ se prononce /k/} & \centering\textbf{La lettre \guillemotleft\ C \guillemotright\ se prononce /s/} \tabularnewline
\hline
\rule{0pt}{1.8em} & \tabularnewline
\hline
\end{tabularx}

\vspace{3pt}
\textbf{Activité 03 (Les enveloppes sonores) :} \textit{Prépare deux enveloppes : une pour le son \textbf{/g/} et l'autre pour \textbf{/3/}.}\\
\textit{Mets chaque mot dans l'enveloppe adéquate :} \textbf{Singe -- gaspiller -- guide -- garçon -- plage -- gomme.}

\vspace{3pt}
\textbf{Activité 04 (Intonation expressive) :} \textit{Lis ces phrases en respectant la ponctuation et l'intonation :}\\
$\bullet$ Les enfants visitent le zoo. \hfill $\bullet$ Aimez-vous les animaux ? \hfill $\bullet$ Quel beau métier ! \hspace{2cm}\null
}

\newpage

% Planche de 15 Vignettes Découpables - Schéma Gazelle
% =========================================================================

\newcommand{\modernGazelleSticker}{%
\begin{tcolorbox}[
  enhanced,
  colback=white,
  colframe=gray!50,
  arc=1.5mm,
  boxrule=0.5pt,
  top=1pt,bottom=1pt,left=1pt,right=1pt,
  height=5.1cm
]
\centering
\includegraphics[width=0.98\linewidth,height=4.8cm,keepaspectratio]{images/gazelle_schema.png}
\end{tcolorbox}
}

\newcommand{\gazelleStickersPage}{%
\noindent
\begin{minipage}[t]{0.315\textwidth}
\modernGazelleSticker
\vspace{-4pt}
\modernGazelleSticker
\vspace{-4pt}
\modernGazelleSticker
\vspace{-4pt}
\modernGazelleSticker
\vspace{-4pt}
\modernGazelleSticker
\end{minipage}
\hfill
\begin{minipage}[t]{0.315\textwidth}
\modernGazelleSticker
\vspace{-4pt}
\modernGazelleSticker
\vspace{-4pt}
\modernGazelleSticker
\vspace{-4pt}
\modernGazelleSticker
\vspace{-4pt}
\modernGazelleSticker
\end{minipage}
\hfill
\begin{minipage}[t]{0.315\textwidth}
\modernGazelleSticker
\vspace{-4pt}
\modernGazelleSticker
\vspace{-4pt}
\modernGazelleSticker
\vspace{-4pt}
\modernGazelleSticker
\vspace{-4pt}
\modernGazelleSticker
\end{minipage}%
}

% Page 10 : Planche Vignettes Gazelle (Exemplaire 1/5)
\gazelleStickersPage


\end{document}
